\section{Zusammenfassung}

LearnHub ist eine Lernplattform die speziell für die 
Maturavorbereitung an der HTL Leonding entwickelt wurde. 
Bestehende Plattformen wie StudyDrive oder Knowunity sind 
auf allgemeinbildende Schulen und Universitäten ausgerichtet 
und bieten keine Unterstützung für HTL-spezifische Fächer 
oder die Struktur der Maturaprüfung. LearnHub schließt 
diese Lücke.

Im Rahmen dieser Diplomarbeit wurden zwei Funktionsbereiche 
entwickelt. Der \textbf{Maturaupload} ermöglicht es 
Schüler:innen eigene Mitschriften hochzuladen und einem 
Fach sowie einem Themenpool zuzuordnen. Lehrkräfte können 
diese Mitschriften prüfen und freigeben, bevor sie für alle 
sichtbar werden. Der \textbf{Fragenkonfigurator} erlaubt 
das Erstellen und Lernen von Fragen, entweder privat für 
die eigene Vorbereitung oder öffentlich für alle 
Nutzer:innen der Plattform.

Technisch läuft das Backend auf Java mit Quarkus, das 
Frontend wurde mit Angular umgesetzt. Die Authentifizierung 
übernimmt Keycloak, der mit dem LDAP-Verzeichnis der HTL 
verbunden ist. Dadurch können Schüler:innen und Lehrkräfte 
ihre gewohnten HTL-Zugangsdaten verwenden, ohne sich extra 
registrieren zu müssen.

Die gesamte Anwendung läuft in Docker-Containern auf einer 
virtuellen Maschine der HTL Leonding. Mit Docker Compose 
und Traefik als Reverse-Proxy wurde eine Infrastruktur 
geschaffen die einfach zu warten und zu aktualisieren ist.

Die Plattform funktioniert stabil und erfüllt alle geplanten 
Anforderungen. Was wir beim nächsten Mal von Anfang an 
anders machen würden: CI/CD. Der manuelle Deployment-Prozess 
hat funktioniert, war aber zeitaufwendig und fehleranfällig. 
Eine GitHub Actions Pipeline wurde zwar bereits aufgesetzt, 
aber eine vollständige Automatisierung von Anfang an hätte 
viel Zeit gespart.

\section{Ausblick}

LearnHub bietet als Lernplattform eine gute Grundlage um künstliche 
Intelligenz sinnvoll einzusetzen. Die bestehende Struktur mit Fächern, 
Themenpools, Fragen und Mitschriften liefert bereits strukturierte 
Daten, auf denen KI-Funktionen aufbauen können.

\subsection{Automatische Fragengenerierung aus Mitschriften}

Eine naheliegende Erweiterung wäre die automatische Generierung von 
Fragen aus hochgeladenen Mitschriften. Beim Upload einer Mitschrift 
könnte ein KI-Modell den PDF-Inhalt analysieren und automatisch 
Fragen mit passenden Antworten vorschlagen. Diese würden dann nicht 
direkt veröffentlicht, sondern zuerst von einer Lehrkraft geprüft 
und freigegeben. Damit würde der Fragenkonfigurator deutlich 
effizienter, da Lehrkräfte nicht jede Frage manuell erstellen 
müssten.

Technisch wäre das über die API eines Large Language Models wie 
GPT-4 oder einem selbst gehosteten Modell umsetzbar. Das Backend 
würde nach dem Upload die PDF-Datei an das Modell schicken und 
die generierten Fragen als Vorschläge in der Datenbank speichern.

\subsection{Intelligente Lernempfehlungen}

Eine weitere Möglichkeit wäre ein personalisiertes Lernsystem. 
Aktuell werden dem Schüler alle Fragen eines Themenpools 
gleichmäßig angezeigt. Mit KI könnte das System analysieren 
welche Fragen ein Schüler häufig falsch beantwortet und diese 
öfter einblenden. Fragen die bereits gut beherrscht werden 
würden seltener erscheinen. Dieses Prinzip nennt sich Spaced 
Repetition und ist aus Apps wie Anki bekannt.

Dafür müsste das Backend Lernfortschritte pro Nutzer und Frage 
speichern. Auf Basis dieser Daten könnte ein einfaches Modell 
berechnen welche Fragen als nächstes gezeigt werden sollen.

\subsection{KI-gestützter Chat-Tutor}

Ähnlich wie bei Knowunity könnte LearnHub einen KI-gestützten 
Chat-Tutor integrieren. Schüler:innen könnten Fragen zu einem 
Themenpool stellen und der Tutor würde auf Basis der vorhandenen 
Mitschriften und Fragen antworten. Der Vorteil gegenüber einem 
allgemeinen KI-Assistenten wäre, dass die Antworten auf die 
HTL-spezifischen Inhalte der Plattform beschränkt wären.

Technisch wäre das über Retrieval-Augmented Generation (RAG) 
umsetzbar. Dabei werden die relevanten Mitschriften aus der 
Datenbank geladen und zusammen mit der Frage des Schülers an 
das Sprachmodell geschickt. Das Modell antwortet dann auf Basis 
dieser konkreten Inhalte.

\subsection{Automatische Bewertung von Freitextantworten}

Aktuell gibt der Lernmodus keine Rückmeldung ob eine Freitextantwort 
richtig oder falsch ist, der Schüler muss das selbst einschätzen. 
Mit KI könnte das System die eingegebene Antwort mit der 
Musterlösung vergleichen und eine automatische Bewertung geben, 
also ob die Antwort inhaltlich korrekt ist, auch wenn sie nicht 
wortwörtlich mit der Musterlösung übereinstimmt.